The two new models solve the inverse problem for the logistic growth law with a
physics-informed neural network (PINN) \citep{raissi2019physics}. Unlike the
classical logistic baseline, the PINN does not evaluate the closed-form
logistic solution during training. Instead, it learns a neural trajectory
$u_\eta(\tau)$ and estimates the physical parameters by penalizing violations
of the differential equation.

The unknowns in the deterministic PINN inverse problem are
$$
  \eta,\quad r,\quad K,
$$
where $\eta$ denotes all neural-network weights and biases. The unknowns in the
stochastic PINN inverse problem are
$$
  \eta,\quad r,\quad K,\quad \sigma.
$$
The role of the neural trajectory $u_\eta$ is to provide a differentiable
surrogate for the unobserved continuous-time biomass curve. The physical
parameters are then identified by forcing this differentiable surrogate to
match the measurements and to satisfy the logistic differential equation at
collocation points between measurements. This is the main distinction from the
closed-form classical logistic fit: the PINN estimates a curve and the
parameters jointly.

\subsubsection{Shared PINN trajectory architecture}

Both logistic PINNs use a scalar neural trajectory
$$
  u_\eta:[0,1]\rightarrow \mathbb R_+,
  \qquad
  \widehat y(t)=u_\eta(\tau(t)).
$$
The network input is the normalized time $\tau$. The hidden architecture has
three fully connected hidden layers, 32 units in each hidden layer, and
hyperbolic tangent activation:
$$
  h_1=\tanh(W_1\tau+b_1),\qquad
  h_\ell=\tanh(W_\ell h_{\ell-1}+b_\ell),\quad \ell=2,3,
$$
followed by one scalar output layer. Positivity of the fitted biomass
trajectory is enforced by
$$
  u_\eta(\tau)=\operatorname{softplus}(W_o h_3+b_o)+10^{-6}.
$$
The trajectory network has 2209 trainable weights and biases:
$$
\begin{aligned}
  \hbox{input to hidden 1:}&\quad 1\cdot 32+32=64,\\
  \hbox{hidden 1 to hidden 2:}&\quad 32\cdot 32+32=1056,\\
  \hbox{hidden 2 to hidden 3:}&\quad 32\cdot 32+32=1056,\\
  \hbox{hidden 3 to output:}&\quad 32\cdot 1+1=33.
\end{aligned}
$$
Thus the deterministic logistic PINN has $2209+2=2211$ fitted parameters after
adding the raw variables for $r$ and $K$, while the stochastic logistic PINN
has $2209+3=2212$ fitted parameters after adding the raw variable for
$\sigma$.

The carrying capacity is constrained to remain above all observed responses,
not only above the first row. Define
$$
  K_{\min}=\max_{(i,j)\in\Omega} y_i^{[j]}+10^{-3}.
$$
The deterministic PINN optimizes unconstrained variables $a,b\in\mathbb R$ and
maps them to physical parameters by
$$
  r=\operatorname{softplus}(a)+10^{-8},
  \qquad
  K=K_{\min}+\operatorname{softplus}(b)+10^{-8}.
$$
The stochastic PINN adds a third unconstrained variable $c$ and uses
$$
  \sigma=\operatorname{softplus}(c)+10^{-8}.
$$
Here $r$ is on the original time scale, because the physics residual includes
the time-scale factor $t_{\max}-t_{\min}=168$.

The PINN initialization is chosen to start from a stable positive trajectory.
The final network layer is initialized so that $u_\eta(\tau)$ is approximately
constant at the mean observed initial level,
$$
  \bar y_0=\frac{1}{J}\sum_{j=1}^J y_0^{[j]}.
$$
The initial carrying capacity is set slightly above the largest observed
response, and the initial growth rate is estimated from the first and last
observed means using the analytic logistic relationship. If that estimate is
not admissible, a small positive default value is used. The stochastic fit
starts from a positive multiplicative diffusion value, $\sigma=0.02$, before
gradient-based optimization.

\subsubsection{Deterministic logistic PINN}

The deterministic inverse problem is based on the original-time logistic ODE
$$
  \frac{dX}{dt}=rX\left(1-\frac{X}{K}\right).
$$
Since the PINN is differentiated with respect to normalized time $\tau$, the
residual uses
$$
  \frac{du_\eta}{d\tau}
  =
  (t_{\max}-t_{\min})\,
  r u_\eta(\tau)\left(1-\frac{u_\eta(\tau)}{K}\right).
$$
Let $\{\xi_m\}_{m=1}^M$ be $M=128$ uniformly spaced collocation points in
$[0,1]$. The data, initial-condition, and physics losses are
$$
  \mathcal L_{\mathrm{data}}(\eta)
  =
  \frac{1}{n}
  \sum_{(i,j)\in\Omega}
  \left(
    \frac{u_\eta(\tau_i)-y_i^{[j]}}{s_y}
  \right)^2,
$$
$$
  \mathcal L_{\mathrm{init}}(\eta)
  =
  \frac{1}{J}
  \sum_{j=1}^J
  \left(
    \frac{u_\eta(0)-y_0^{[j]}}{s_y}
  \right)^2,
$$
and
$$
  \mathcal L_{\mathrm{phys}}(\eta,r,K)
  =
  \frac{1}{M}
  \sum_{m=1}^M
  \left[
    \frac{
      \dfrac{du_\eta}{d\tau}(\xi_m)
      -
      (t_{\max}-t_{\min})r u_\eta(\xi_m)
      \left(1-\dfrac{u_\eta(\xi_m)}{K}\right)
    }{s_y}
  \right]^2 .
$$
The deterministic PINN objective is
$$
  \mathcal L_{\mathrm{detPINN}}
  =
  \mathcal L_{\mathrm{data}}
  +
  \lambda_0\mathcal L_{\mathrm{init}}
  +
  \lambda_{\mathrm{phys}}\mathcal L_{\mathrm{phys}},
$$
with $\lambda_0=1$ and $\lambda_{\mathrm{phys}}=1$ in the reported run. This
fit uses one population trajectory against all three replicate trials. It is
therefore more flexible than the classical two-parameter logistic regression,
because the neural trajectory can absorb deviations from the exact analytic
curve, but it is still regularized toward the logistic differential equation.
Training used Adam for 3000 epochs with learning rate $5\times10^{-3}$.

The deterministic PINN inverse problem is solved by the following optimization
loop:
\begin{enumerate}
  \item Build the observation tensor, mask missing entries, compute $s_y$, and
        map observation times to $\tau_i\in[0,1]$.
  \item Create the collocation grid
        $\xi_1,\ldots,\xi_M$ on $[0,1]$. These points are not additional data;
        they are locations where the logistic differential equation is
        enforced.
  \item Initialize the trajectory network $u_\eta$ and raw physical variables
        $a,b$. The reported parameters at every epoch are obtained from
        $r=\operatorname{softplus}(a)+10^{-8}$ and
        $K=K_{\min}+\operatorname{softplus}(b)+10^{-8}$.
  \item Evaluate the population trajectory at the observation times to compute
        $\mathcal L_{\mathrm{data}}$ and $\mathcal L_{\mathrm{init}}$.
  \item Evaluate $u_\eta(\xi_m)$ at the collocation points and use automatic
        differentiation to compute $du_\eta/d\tau$ at those same points.
  \item Form the residual
        $du_\eta/d\tau-(t_{\max}-t_{\min})r u_\eta(1-u_\eta/K)$ and compute
        $\mathcal L_{\mathrm{phys}}$.
  \item Add the weighted terms to obtain
        $\mathcal L_{\mathrm{detPINN}}$, backpropagate through both the
        network and the softplus-transformed parameters, and update
        $(\eta,a,b)$ with Adam.
  \item After training, transform the final raw variables to
        $\widehat r$ and $\widehat K$, and evaluate $u_{\widehat\eta}(\tau)$
        on the observed and dense plotting grids.
\end{enumerate}

The gradients with respect to $r$ and $K$ are not computed by finite
differences. They are obtained by automatic differentiation through the loss
terms. In particular, the physics loss connects $r$ and $K$ directly to the
neural trajectory derivative, so the fitted parameters are those that make the
learned curve both close to the observations and dynamically consistent with
logistic growth.